\section{Conclusion}
\label{sec:conclusion}

We revisited anisotropy in large language model training and identified a sharper structural principle:
\emph{activation anisotropy is predominantly driven by a coherent rank-one mean bias.}
This mean component arises naturally from hierarchical semantic aggregation
and accumulates through residual pathways,
systematically regenerating coherent directions in representation space.

In high-dimensional settings, even modest coordinate-wise bias scales as
$\|\mu\|_2 \sim \sqrt{H}\,\bar{\mu}$,
amplifying small drift into dominant elementwise extremes.
Because blockwise low-bit quantization scales are determined by these extremes,
mean bias becomes the principal driver of dynamic-range inflation and numerical instability.
We show that this dominant rank-one component can be removed through a simple
source-level mean subtraction,
providing effective spectral conditioning using only reduction and elementwise operations.

These findings recast anisotropy as both a \emph{curse} and a \emph{blessing}:
the same structured mean bias that destabilizes low-bit training
also exposes a low-dimensional handle for efficient mitigation.
By isolating and removing this component at its source,
we obtain a hardware-friendly path to stable low-bit LLM training.